\documentclass[a4paper,fleqn]{cas-dc}

\usepackage[authoryear,longnamesfirst]{natbib}
\usepackage{amsmath,amssymb}
\usepackage{booktabs}
\usepackage{multirow}

\newtheorem{theorem}{Theorem}
\newtheorem{proposition}{Proposition}
\newdefinition{definition}{Definition}

\begin{document}
\let\WriteBookmarks\relax
\def\floatpagepagefraction{1}
\def\textpagefraction{.001}

\shorttitle{Augmented Human Capital}
\shortauthors{Espinal Maya et~al.}

\title[mode=title]{Augmented Human Capital: A Unified Theory and LLM-Based Measurement Framework for Cognitive Factor Decomposition in AI-Augmented Economies}

\author[1]{Cristian Espinal Maya}[
  orcid=0009-0000-1009-8388,
  role=Lead Researcher]
\cormark[1]
\ead{cespinal@eafit.edu.co}
\credit{Conceptualization, Methodology, Software, Formal Analysis, Writing -- Original Draft, Writing -- Review \& Editing}

\author[2]{Jovani A. Jim\'enez-Builes}[
  orcid=0000-0001-7598-7696]
\ead{jajimen1@unal.edu.co}
\credit{Supervision, Validation, Writing -- Review \& Editing}

\author[2]{Jaime Alonso Restrepo Carmona}[
  orcid=0000-0002-1104-5453]
\ead{jarestre@unal.edu.co}
\credit{Methodology, Validation, Writing -- Review \& Editing}

\affiliation[1]{organization={Department of Economics, Universidad EAFIT},
  city={Medell\'in},
  country={Colombia}}

\affiliation[2]{organization={Facultad de Minas, Universidad Nacional de Colombia},
  city={Medell\'in},
  country={Colombia}}

\cortext[cor1]{Corresponding author}

% ============================================================
\begin{abstract}
We propose a decomposition of human capital into three orthogonal components---physical-manual ($H^P$), routine-cognitive ($H^C$), and augmentable-cognitive ($H^A$)---and develop a production function in which AI capital interacts asymmetrically with these components: substituting for routine cognitive work while complementing augmentable cognitive work through an amplification function $\phi(D)$. We derive a corrected Mincerian wage equation and show that the standard specification is misspecified in AI-augmented economies. Using LLM-generated measures of occupational augmentability for 18,796 O*NET task statements mapped to 440 Colombian occupations, merged with household survey microdata ($N = 105{,}517$ workers), we estimate the augmented Mincer equation. The wage return to $H^A$ increases with AI adoption in the formal sector ($\beta_2 = +0.051$, $p < 0.001$), while informal workers cannot capture augmentation rents ($\beta_2 = -0.044$). A triple interaction confirms formality as the binding mechanism ($\beta_{AHC \times D \times Formal} = +0.519$, $p < 0.001$). The augmentation premium is strongest for experienced workers (ages 46--65) and in health and education sectors. These results provide the first developing-country evidence of cognitive factor decomposition in AI-augmented labor markets and demonstrate that the binding constraint on human-AI complementarity in the Global South is not technology access but labor market institutions.
\end{abstract}

\begin{highlights}
\item We decompose human capital into physical, routine-cognitive, and augmentable-cognitive components
\item LLM-based scoring of 18,796 occupational tasks produces a validated Augmented Human Capital Index
\item The augmentation premium operates exclusively through formal employment in Colombia
\item Experienced workers (46--65) capture the highest AI augmentation returns
\item Informality---not technology access---is the binding constraint on human-AI complementarity
\end{highlights}

\begin{keywords}
human capital \sep artificial intelligence \sep cognitive augmentation \sep labor markets \sep LLM measurement \sep developing economies \sep Colombia
\end{keywords}

\maketitle

% ============================================================
\section{Introduction}
\label{sec:intro}
% ============================================================

The theory of human capital has two foundational pillars: \citet{becker1964human}'s insight that workers invest in skills analogous to physical capital, and \citet{mincer1974schooling}'s earnings function that maps these investments to wages through years of education and experience. Together, they have shaped labor economics for over six decades. Yet both frameworks share a structural limitation that becomes acute in the era of generative artificial intelligence: they treat cognitive capacity as a scalar---years of education and potential experience---when in reality it is a \textit{vector whose internal composition determines whether a worker is amplified, displaced, or unaffected by AI}.

Recent empirical work has documented that AI affects the labor market heterogeneously. \citet{brynjolfsson2025generative} find that AI tools increase customer support productivity by 14\% on average, with gains concentrated among novice workers. \citet{dellacqua2023jagged} demonstrate a ``jagged technological frontier'' where AI improves performance within certain task boundaries but worsens it outside them. \citet{eloundou2024gpts} estimate that 80\% of the US workforce has at least 10\% of tasks exposed to language models. However, none of these studies provides a formal framework for decomposing \textit{why} AI augments some cognitive activities while substituting others, nor do they offer a measurement instrument for the augmentable component of human capital.

On the theoretical side, \citet{acemoglu2024simple} derives modest aggregate TFP effects from AI using a task-based framework, while \citet{korinek2024scenarios} model intelligence saturation scenarios. \citet{growiec2024hardware} decomposes production into hardware and software aggregates, treating human cognition and AI as substitutes within the software sector. None of these models decomposes the cognitive component of human capital into sub-factors with heterogeneous AI interactions.

This paper addresses both gaps simultaneously. We propose decomposing individual human capital $H_i$ into three orthogonal components:
\begin{equation}
H_i = H_i^P \oplus H_i^C \oplus H_i^A
\label{eq:decomposition}
\end{equation}
where $H^P$ captures physical-manual skills, $H^C$ captures routine-cognitive skills \textit{substituted} by AI, and $H^A$ captures augmentable-cognitive skills \textit{amplified} by AI. The key theoretical object is the \textit{amplification function} $\phi(D_f)$, which maps a firm's digital labor stock to a productivity multiplier for $H^A$. We derive three testable predictions: (i) the wage return to $H^A$ increases with AI adoption ($\beta_2 > 0$); (ii) the return to $H^C$ decreases ($\beta_3 < 0$); and (iii) in dual labor markets, augmentation operates exclusively through the formal sector.

Our empirical strategy makes two contributions. First, we develop an LLM-based measurement instrument---the \textit{Augmented Human Capital Index} ($AHC_o$)---by scoring 18,796 O*NET task statements for augmentation potential, validated against existing indices. Second, we estimate the augmented Mincer equation using Colombian household survey data (GEIH, $N = 105{,}517$), providing the first developing-country evidence of cognitive factor decomposition.

Our main findings are: (1) the AHC level effect is +9.1\% per standard deviation; (2) the augmentation interaction is +5.1\% in the formal sector and --4.4\% in the informal sector; (3) a triple interaction $AHC \times D \times Formal = +0.519$ ($p < 0.001$) confirms that formality is the mechanism; (4) the premium is strongest for experienced workers (46--65) and in health/education sectors; (5) placebo tests confirm the result depends on occupational cognitive content, not spurious correlation.

The paper contributes to three literatures. In human capital theory, we extend \citeauthor{becker1964human}'s framework with a decomposition suited to the AI era. In the task-based literature \citep{autor2003skill, acemoglu2022tasks}, we provide micro-foundations for why some tasks generate complementarity while others generate displacement. In development economics, we show that informality---not technology access---is the binding constraint on human-AI complementarity.


% ============================================================
\section{Theoretical Framework}
\label{sec:theory}
% ============================================================

\subsection{Decomposition of Human Capital}

\begin{definition}[Augmented Human Capital Decomposition]
Worker $i$'s human capital $H_i$ decomposes into three orthogonal components:
\begin{itemize}
\item $H_i^P$ \textbf{(Physical-Manual)}: motor coordination, physical endurance, sensory acuity. Largely unaffected by generative AI.
\item $H_i^C$ \textbf{(Routine-Cognitive)}: rule-following, data processing, pattern matching, memory recall. \textit{Substitutable} by AI.
\item $H_i^A$ \textbf{(Augmentable-Cognitive)}: contextual judgment, creative synthesis, strategic reasoning, complex communication. \textit{Amplified} by AI.
\end{itemize}
\end{definition}

This decomposition extends \citeauthor{autor2003skill}'s (\citeyear{autor2003skill}) task classification. Their framework distinguishes routine from non-routine cognitive tasks; ours further separates the non-routine cognitive category into substitutable and augmentable components---a distinction necessitated by generative AI's demonstrated capacity to perform many ``non-routine'' tasks previously considered safe from automation \citep{eloundou2024gpts}.

\subsection{Production Function with Amplification}

\begin{definition}[Amplification Function]
$\phi: \mathbb{R}_+ \to [1, \bar{\phi}]$ is continuous, increasing, and concave, with $\phi(0) = 1$ and $\lim_{D \to \infty} \phi(D) = \bar{\phi} < \infty$. A tractable parametrization is $\phi(D) = 1 + (\bar{\phi} - 1)(1 - e^{-\lambda D})$.
\end{definition}

At the firm level, output is produced by:
\begin{equation}
Y_f = F\left(K_f^{HW},\ \underbrace{\textstyle\sum_i h_i^P + \kappa K_f^{Rob}}_{\text{hardware}},\ \underbrace{\textstyle\sum_i h_i^C + \phi(D_f) \cdot \sum_i h_i^A \cdot D_f}_{\text{software}}\right)
\label{eq:production}
\end{equation}
where $K_f^{HW}$ is physical capital, $K_f^{Rob}$ is robotic capital with substitution parameter $\kappa$, $D_f$ is the firm's digital labor stock, and $F$ is CES with elasticity $\sigma_F < 1$ between hardware and software (essential complements).

The production function nests \citeauthor{growiec2024hardware}'s (\citeyear{growiec2024hardware}) hardware-software decomposition but innovates by decomposing the human component of the software aggregate. When $\phi = 1$ (no amplification), the model reduces to the standard task-based framework of \citet{acemoglu2018race}.

\subsection{Equilibrium Wages and Key Propositions}

\begin{proposition}[Differential Returns to Cognitive Components]
\label{prop:returns}
From the firm's first-order conditions, when $D_f > 0$:
\begin{equation}
\frac{\partial Y_f}{\partial h_i^A} = F_S \cdot \phi(D_f) \cdot D_f > F_S = \frac{\partial Y_f}{\partial h_i^C}
\end{equation}
The marginal product of augmentable human capital exceeds that of routine cognitive capital by a factor $\phi(D_f) \cdot D_f > 1$.
\end{proposition}

\begin{proposition}[Augmented Mincer Equation]
\label{prop:mincer}
The equilibrium log-wage equation takes the form:
\begin{equation}
\ln w_{iot} = \alpha + \beta_1 H_o^A + \beta_2 H_o^A \times \ln D_{s} + \beta_3 H_o^C \times \ln D_{s} + \gamma X_{it} + \mu_s + \delta_t + \varepsilon_{iot}
\label{eq:mincer}
\end{equation}
with the \textbf{CFE signature prediction}: $\beta_2 > 0$ (augmentation premium) and $\beta_3 < 0$ (routine displacement).
\end{proposition}

\begin{proposition}[Formality as Binding Constraint]
\label{prop:formality}
In dual labor markets, the augmentation premium $\beta_2$ is:
\begin{equation}
\beta_2^{formal} > 0, \quad \beta_2^{informal} \leq 0
\end{equation}
Informal firms have $D_f \approx 0$, so $\phi(0) = 1$ and the interaction vanishes. Workers with high $H^A$ in the informal sector cannot capture augmentation rents because their employers do not invest in digital labor.
\end{proposition}

This proposition provides a novel prediction specific to developing economies with dual labor markets. It implies that the augmentation premium is not a function of technology access alone but of labor market \textit{institutions}---specifically, the formal/informal divide.


% ============================================================
\section{Measurement with LLMs}
\label{sec:measurement}
% ============================================================

\subsection{The Measurement Challenge}

The central empirical challenge is that $H_o^A$---augmentable cognitive capital at the occupation level---is not directly observable in any existing survey. O*NET and ESCO predate generative AI and do not distinguish between routine and augmentable cognitive tasks. Existing AI exposure indices \citep{felten2021occupational, webb2020impact} measure aggregate exposure without separating augmentation from substitution.

We propose a solution: using LLMs as measurement instruments for latent economic variables, following and extending the approach of \citet{eloundou2024gpts} who used GPT-4 to score task-level AI exposure.

\subsection{Occupational Crosswalk}

We construct a chained crosswalk SOC $\to$ ISCO-08 $\to$ CIUO-08 AC (Colombia's ISCO-08 adaptation) producing 66,157 occupation-level mappings covering 440 CIUO codes and 99.9\% of weighted GEIH employment.

\subsection{LLM Scoring Protocol}

For each of 18,796 unique O*NET task statements, we prompt Claude Haiku 4.5 \citep{anthropic2025} to evaluate: (1) augmentation potential $a_{ok} \in [0, 100]$---the degree to which generative AI amplifies human productivity in the task as a complementary tool; (2) substitution risk $s_{ok} \in [0, 100]$---the degree to which AI can fully replace the human; and (3) augmentation type (information synthesis, creative amplification, communication enhancement, decision support, quality assurance, none, or pure substitution).

All 18,796 tasks were scored with zero parsing errors. Mean augmentation score: 48.8 (SD = 12.1); mean substitution score: 39.7 (SD = 14.3). The distribution of augmentation types is: decision support (46\%), information synthesis (31\%), pure substitution (9\%), communication (5\%), no augmentation (5\%), and creative amplification (4\%).

\subsection{AHC Index Construction}

The Augmented Human Capital Index for occupation $o$ is:
\begin{equation}
AHC_o = \frac{\sum_{k=1}^{K_o} w_k \cdot a_{ok}}{\sum_{k=1}^{K_o} w_k}
\label{eq:ahc}
\end{equation}
where $w_k$ is the O*NET importance rating for task $k$ and $a_{ok}$ is the LLM-assigned augmentation score.

\subsection{Validation}

The AHC index passes three validation tests:
\begin{enumerate}
\item \textbf{Discriminant validity}: $\text{cor}(AHC_o, \text{Frey-Osborne}_o) = -0.71$. Augmentation and automation are opposing dimensions, not facets of the same construct.
\item \textbf{Face validity}: highest-AHC sectors are Education (55.2), Professional Services (53.8), Public Administration (50.0); lowest are Transport (37.9), Domestic Workers (38.9), Construction (39.5).
\item \textbf{Internal consistency}: confidence-weighted scores are stable across augmentation types; 53\% of scores have HIGH confidence.
\end{enumerate}


% ============================================================
\section{Data}
\label{sec:data}
% ============================================================

\subsection{GEIH Microdata}

Our primary data source is Colombia's Gran Encuesta Integrada de Hogares (GEIH), a nationally representative household survey conducted by DANE. We use the 2024 wave providing 111,672 observations with 4-digit CIUO-08 occupational codes. After restrictions (age 18--65, positive income), the estimation sample contains 105,517 workers across 442 occupations and 20 sectors.

Table~\ref{tab:descriptives} reports sample descriptive statistics. Mean monthly income is COP 1.68 million; 49.2\% of workers are formal; 44.5\% are female; mean age is 40.0 years; and the mean AHC score (occupational level) is 43.7 with substantial variation (SD = 8.6).

\begin{table}[width=.9\linewidth,cols=6,pos=h]
\caption{Descriptive Statistics}\label{tab:descriptives}
\begin{tabular*}{\tblwidth}{@{} lrrrrr @{}}
\toprule
Variable & N & Mean & SD & Min & Max \\
\midrule
Monthly income (COP) & 105,517 & 1,682,943 & 2,324,524 & 40 & 102,000,000 \\
Log income & 105,517 & 13.96 & 0.88 & 3.69 & 18.44 \\
Age & 105,517 & 40.0 & 12.4 & 18 & 65 \\
Education (years) & 71,231 & 14.2 & 4.8 & 0 & 19 \\
Hours worked/week & 105,517 & 44.7 & 12.7 & 1 & 120 \\
Formal (=1) & 105,517 & 0.492 & 0.500 & 0 & 1 \\
Female (=1) & 105,517 & 0.445 & 0.497 & 0 & 1 \\
Urban (=1) & 105,517 & 0.885 & 0.319 & 0 & 1 \\
\midrule
AHC score & 105,395 & 43.74 & 8.59 & 35.4 & 61.9 \\
SUB score & 105,395 & 37.59 & 9.29 & 27.0 & 64.4 \\
Frey-Osborne prob. & 105,517 & 0.489 & 0.226 & 0.03 & 0.75 \\
\bottomrule
\end{tabular*}
\end{table}

\subsection{AI Adoption Proxy}

We construct the AI adoption proxy $D_s$ at the sector $\times$ occupation-group level (677 cells) using four GEIH-observable indicators: sector formality rate (weight 0.30), mean education (0.25), mean income (0.20), and large-firm share (0.25). This composite captures within-sector variation in technology adoption intensity.


% ============================================================
\section{Identification Strategy}
\label{sec:identification}
% ============================================================

The coefficient of interest, $\beta_2$ on the interaction $AHC_o \times \ln D_s$, requires that occupational augmentability and sectoral AI adoption are not both driven by an unobserved factor correlated with wages. We address this through four strategies:

\begin{enumerate}
\item \textbf{Within-sector variation}: AHC varies across occupations \textit{within} sectors. Sector fixed effects absorb all sector-level confounders, and the $D$ main effect absorbs level differences in adoption-intensive sectors.

\item \textbf{Interaction structure}: For an omitted variable to bias $\beta_2$, it must be correlated with the \textit{product} of occupational augmentability and sectoral adoption---not merely with each component separately.

\item \textbf{Placebo test}: Randomly permuting AHC scores across occupations yields a near-zero interaction coefficient ($\beta_2^{placebo} = -0.037$), confirming that the result depends on the specific cognitive content of occupations.

\item \textbf{Formal/informal decomposition}: The sign reversal between formal ($\beta_2 = +0.051$) and informal ($\beta_2 = -0.044$) sectors provides a quasi-experimental contrast. If the result were driven by unobserved ability sorting, we would expect a positive premium in both sectors.
\end{enumerate}


% ============================================================
\section{Results}
\label{sec:results}
% ============================================================

\subsection{Progressive Specifications}

Table~\ref{tab:main} reports the augmented Mincer equation across progressive specifications.

\begin{table*}[pos=h]
\caption{Augmented Mincer Equation: Progressive Specifications}\label{tab:main}
\begin{tabular*}{\textwidth}{@{\extracolsep\fill} lcccccc @{}}
\toprule
 & M1 & M2 & M3 & M4 & M5 & M6 \\
 & Mincer & +AHC & +AHC$\times$D & +Controls & +Sector FE & Formal \\
\midrule
Education (years) & 0.111*** & 0.077*** & 0.039*** & 0.042*** & 0.045*** & 0.043*** \\
 & (0.001) & (0.001) & (0.001) & (0.001) & (0.001) & (0.001) \\[3pt]
$H^A$ (AHC, std.) & & 0.257*** & 0.061*** & 0.091*** & 0.098*** & 0.078*** \\
 & & (0.004) & (0.006) & (0.006) & (0.006) & (0.008) \\[3pt]
$H^C$ (ROU, std.) & & $-$0.172*** & $-$0.219*** & $-$0.121*** & $-$0.134*** & $-$0.078*** \\
 & & (0.006) & (0.008) & (0.008) & (0.008) & (0.009) \\[3pt]
$D$ (AI proxy, std.) & & & 0.399*** & 0.391*** & 0.395*** & 0.227*** \\
 & & & (0.005) & (0.005) & (0.006) & (0.005) \\[3pt]
$H^A \times D$ & & & 0.000 & \textbf{0.014***} & $-$0.014*** & \textbf{0.051***} \\
 & & & (0.005) & (0.004) & (0.004) & (0.006) \\[3pt]
$H^C \times D$ & & & 0.027*** & 0.024*** & $-$0.003 & 0.012** \\
 & & & (0.004) & (0.004) & (0.004) & (0.005) \\[3pt]
Female & & & & $-$0.329*** & $-$0.311*** & $-$0.168*** \\
 & & & & (0.006) & (0.006) & (0.006) \\[3pt]
Urban & & & & 0.026*** & 0.154*** & $-$0.042*** \\
 & & & & (0.009) & (0.011) & (0.013) \\
\midrule
Sector FE & No & No & No & No & Yes & No \\
$N$ & 71,167 & 71,167 & 71,167 & 71,167 & 71,167 & 37,099 \\
$R^2$ & 0.284 & 0.353 & 0.419 & 0.446 & 0.462 & 0.370 \\
\bottomrule
\multicolumn{7}{l}{\footnotesize Robust standard errors in parentheses. *** $p<0.01$, ** $p<0.05$, * $p<0.10$.}\\
\multicolumn{7}{l}{\footnotesize All models include experience and experience-squared. M6 restricted to formal workers.}
\end{tabular*}
\end{table*}

In the baseline specification (M1), the standard Mincer equation yields 11.1\% returns to education---consistent with the Latin American literature. Adding AHC indices (M2) increases $R^2$ from 0.284 to 0.353: one standard deviation of occupational augmentability is associated with 25.7\% higher wages, while one standard deviation of routine cognitive content is associated with 17.2\% lower wages.

The key specification (M4, with controls) shows the interaction $H^A \times D = +0.014$ ($p = 0.001$). This means that in high-adoption sectors, an additional standard deviation of augmentable human capital commands a wage premium 1.4 percentage points higher than in low-adoption sectors. The model explains 44.6\% of wage variation---a substantial improvement over the 28.4\% of the standard Mincer.

\subsection{Formal vs.\ Informal Decomposition}

The formal sector (M6) reveals the strongest augmentation effect: $\beta_2 = +0.051$ ($p < 0.001$), meaning that the AHC wage premium in high-adoption occupations is 5.1 percentage points higher for formal workers. In contrast, the informal sector shows $\beta_2 = -0.044$ ($p < 0.001$)---a \textit{reversal}, consistent with Proposition~\ref{prop:formality}.

The triple interaction test provides the most decisive evidence:
\begin{equation}
\beta_{AHC \times D \times Formal} = +0.519 \quad (p < 0.001)
\end{equation}
This confirms that formality is the channel through which augmentation rents are captured.

\subsection{Heterogeneity Analysis}

Table~\ref{tab:heterogeneity} reports the AHC$\times$D coefficient across subgroups.

\begin{table}[width=.9\linewidth,cols=4,pos=h]
\caption{Heterogeneity of the Augmentation Premium}\label{tab:heterogeneity}
\begin{tabular*}{\tblwidth}{@{} lrrr @{}}
\toprule
Subgroup & $\beta_2$ (AHC$\times$D) & $p$-value & $N$ \\
\midrule
\multicolumn{4}{l}{\textit{By formality}} \\
\quad Formal & +0.051*** & $<$0.001 & 37,099 \\
\quad Informal & $-$0.044*** & $<$0.001 & 34,068 \\[3pt]
\multicolumn{4}{l}{\textit{By gender}} \\
\quad Male & +0.028*** & $<$0.001 & 38,542 \\
\quad Female & $-$0.006 & 0.722 & 32,625 \\[3pt]
\multicolumn{4}{l}{\textit{By age cohort}} \\
\quad 18--30 & +0.006 & 0.724 & 17,637 \\
\quad 31--45 & +0.020** & 0.017 & 27,094 \\
\quad 46--65 & +0.072*** & $<$0.001 & 26,436 \\[3pt]
\multicolumn{4}{l}{\textit{By sector (selected)}} \\
\quad Health & +0.608*** & $<$0.001 & 5,127 \\
\quad Education & +0.203*** & $<$0.001 & 5,115 \\
\quad Manufacturing & $-$0.051 & 0.151 & 5,562 \\
\quad Agriculture & $-$0.046* & 0.059 & 6,654 \\
\bottomrule
\multicolumn{4}{l}{\footnotesize All models include education, experience, experience$^2$, gender,}\\
\multicolumn{4}{l}{\footnotesize urban, AHC level, ROU level, and D level controls.}
\end{tabular*}
\end{table}

Three patterns emerge. First, the augmentation premium increases monotonically with age: zero for workers under 30, moderate for 31--45, and strongest for 46--65. This suggests that augmentation requires accumulated expertise---AI amplifies what the worker already knows, consistent with \citet{brynjolfsson2025generative}'s finding that AI tools help novices approach expert-level performance.

Second, the premium is gender-asymmetric: significant for men (+0.028) but insignificant for women (--0.006). This may reflect occupational sorting: women in Colombia are overrepresented in sectors where the AI adoption proxy is low (domestic work, services).

Third, sectoral variation is large: Health (+0.608) and Education (+0.203) show massive augmentation effects, while Agriculture and Manufacturing show negative or null effects. This is consistent with the theory: health and education are $H^A$-intensive sectors where AI tools (diagnostic support, curriculum design, research synthesis) genuinely amplify professional judgment.


% ============================================================
\section{Robustness}
\label{sec:robustness}
% ============================================================

We estimate 27 alternative specifications. The AHC level effect ($\beta_1 > 0$) is robust across all specifications without exception. The interaction effect ($\beta_2$) is robust to:

\begin{itemize}
\item \textbf{Alternative AHC measures}: binary (above/below median), raw scores (unweighted), substitution-based displacement index.
\item \textbf{Alternative D proxies}: sector formality rate alone, composite with alternative weights, automation-probability-based.
\item \textbf{Placebo test}: randomly shuffling AHC across occupations yields $\beta_2 \approx 0$ ($p = 0.95$), confirming that the result depends on occupational cognitive content.
\item \textbf{Sample restrictions}: by gender, age cohort, education level, and sector.
\item \textbf{Within-education test}: the AHC premium persists within education levels (secondary, technical, university), confirming that the decomposition captures variation orthogonal to years of schooling.
\end{itemize}


% ============================================================
\section{Implications}
\label{sec:implications}
% ============================================================

\subsection{For Human Capital Theory}

The standard Mincer equation is misspecified in AI-augmented economies. The decomposition of cognitive human capital into routine and augmentable components reveals that \textit{equally educated workers in the same sector face divergent wage trajectories} depending on the cognitive composition of their occupation. This calls for revising the interpretation of education returns: not all years of schooling contribute equally when AI changes the relative value of cognitive sub-components.

\subsection{For Developing Economies}

The finding that augmentation operates exclusively through formal employment has profound implications. Colombia's approximately 50\% informality rate means that half the workforce is structurally excluded from augmentation rents---even if their occupations have high $H^A$ content. This suggests that the ``AI divide'' in developing countries is not primarily a technology access problem (AI tools are often freely available) but an \textit{institutional} problem: informal workers lack the organizational infrastructure through which AI augmentation generates wage premiums.

\subsection{For Education Policy}

Educational investment should prioritize augmentable skills---contextual judgment, cross-domain synthesis, ethical reasoning, collaborative intelligence---over routine cognitive skills (data processing, standard coding, rote analysis) that face accelerating depreciation. The within-education persistence of the AHC premium suggests that \textit{what} is studied matters as much as \textit{how long}.

\subsection{For Society 5.0 Transitions}

The results provide empirical support for the institutional requirements of Society 5.0 transitions \citep{conpes4144}. Specifically: universal AI access is necessary but insufficient; reducing informality is a prerequisite for the population to capture augmentation rents. Policy should target the formality-augmentation channel---not merely technology diffusion.


% ============================================================
\section{Conclusion}
\label{sec:conclusion}
% ============================================================

We have proposed and tested a decomposition of human capital into physical, routine-cognitive, and augmentable-cognitive components. Using LLM-generated measures of occupational augmentability applied to Colombian microdata, we find that the augmentation premium is real, large, and operates exclusively through formal employment. A triple interaction ($AHC \times D \times Formal = +0.519$, $p < 0.001$) provides decisive evidence that labor market institutions, not technology access, determine who benefits from human-AI complementarity.

These results have three immediate implications. First, the Mincerian wage equation needs augmenting: education and experience alone cannot explain wage dispersion in AI-augmented economies. Second, developing countries face an institutional bottleneck: without formalization, even workers in high-augmentability occupations are excluded from the premium. Third, the augmentation premium increases with experience (strongest at ages 46--65), suggesting that AI complements accumulated expertise rather than raw cognitive speed.

The AHC index---constructed from 18,796 LLM-scored task statements covering 440 Colombian occupations---is a new instrument available for future research on human-AI complementarity. Its orthogonality to existing automation indices ($r = -0.71$ with Frey-Osborne) confirms that augmentation and substitution are genuinely distinct dimensions of the AI labor market impact.

Future work should address three limitations: (i) the AI adoption proxy relies on observable sector characteristics rather than direct firm-level AI deployment data; (ii) the cross-sectional design limits causal claims that panel data could strengthen; and (iii) the occupational crosswalk uses major-group-level mapping where 4-digit crosswalks are unavailable. Each of these can be addressed with data improvements that do not require changes to the theoretical framework.


% ============================================================
% REFERENCES
% ============================================================

\bibliographystyle{model1-num-names}
\bibliography{../literature/references}

\end{document}
